\section{Trabalhos Relacionados}\label{sec:related}
Diversas pesquisas têm sido desenvolvidas para verificar o nível de atenção dos alunos em ambientes educacionais mediados por tecnologia, especialmente impulsionadas pela pandemia de COVID-19, que levou à adoção do Ensino Remoto Emergencial (ERE). Mesmo não sendo mais emergencial, a modalidade remota ganhou grande popularidade após esse período, com um número crescente de alunos matriculados, como identificado no estudo de \cite{moran2009ensino}.

Nesse cenário, Ong et al. \cite{ong2021application} propuseram a detecção de atenção com base na captura de imagens dos alunos. Para distinguir alunos focados daqueles dispersos, utilizaram-se os algoritmos Viola-Jones e Sobel Edge, com o objetivo de identificar componentes faciais (como rosto, boca e olhos) e detectar o estado dos olhos (abertos ou fechados), respectivamente. Justificam o uso de imagens — um tipo de dado extremamente sensível — com base em dois grandes desafios enfrentados na detecção de atenção. O primeiro refere-se à alternativa de solicitar que os alunos interajam respondendo a perguntas; no entanto, essa abordagem pode ser ineficaz, já que nem todos respondem e, quando o fazem, podem apresentar lentidão, o que atrasa o andamento da aula. A segunda alternativa seria solicitar que os alunos mantenham a câmera ligada, mas, além de exigir a verificação manual por parte do professor, os alunos ainda poderiam utilizar outros dispositivos, relegando a aula a segundo plano. Apesar de o estudo ter obtido sucesso na identificação dos elementos faciais, foi observado que houve falhas na detecção consistente do estado dos olhos ao aplicar o algoritmo Sobel Edge.

Outros trabalhos, como o de Sasha et al. (2021) \cite{saha2021attendance}, constataram, por meio de análise qualitativa, que 80\% dos alunos entrevistados relataram déficit de atenção durante as aulas remotas, enquanto 40\% relataram que compareciam apenas para registrar presença, sem engajamento ou atenção efetiva. Diante disso, esta pesquisa propõe implementar um software monitor que utiliza dados de imagem dos usuários combinados com modelos de inteligência artificial (IA) para avaliar o grau de atenção do indivíduo e, com base nisso, atribuir a presença. O programa monitoraria a conectividade de rede dos alunos, ativaria as câmeras e, em intervalos regulares, capturaria imagens dos usuários. O sistema processaria essas imagens localmente no cliente, verificando por quanto tempo o aluno esteve olhando para fora da tela e comparando os resultados com modelos treinados com IA. Por fim, os professores apenas utilizariam os resultados fornecidos pelo sistema, referentes ao nível de atenção de cada aluno.

Diferentemente desses trabalhos, o presente estudo propõe uma abordagem menos invasiva e mais centrada na experiência do aluno. Por meio da coleta de dados de navegação, a ferramenta Engajamento e Supervisão do Processo Educacional Online (ESPEON) busca identificar padrões de atenção e engajamento sem recorrer a imagens ou gravações. Esses indicadores servirão como base para uma análise crítica da eficácia da metodologia expositiva em ambientes remotos de ensino, respeitando os princípios de privacidade e a conformidade com a Lei Geral de Proteção de Dados (LGPD).
